\documentclass[12pt,a4paper]{article}

% =====================================================
% Mathematics
% =====================================================
\usepackage{amsmath,amssymb,amsfonts}
\usepackage{mathtools}

% =====================================================
% Graphics
% =====================================================
\usepackage{graphicx}
\usepackage{float}
\usepackage{subcaption}

% =====================================================
% Tables
% =====================================================
\usepackage{booktabs}
\usepackage{array}
\usepackage{multirow}
\usepackage{longtable}
\usepackage{tabularx}
\usepackage{makecell}
\usepackage{adjustbox}

% =====================================================
% Lists
% =====================================================
\usepackage{enumitem}

% =====================================================
% Page Layout
% =====================================================
\usepackage[a4paper,margin=1in]{geometry}

% =====================================================
% Hyperlinks
% =====================================================
\usepackage[hidelinks]{hyperref}

% =====================================================
% Captions
% =====================================================
\usepackage{caption}

% =====================================================
% Typography
% =====================================================
\usepackage{microtype}
\usepackage{setspace}
\onehalfspacing

% =====================================================
% Paragraph Formatting
% =====================================================
\usepackage{parskip}

% =====================================================
% Colors
% =====================================================
\usepackage{xcolor}

% =====================================================
% Float Control
% =====================================================
\usepackage{placeins}

% =====================================================
% Landscape Pages (Wide Tables)
% =====================================================
\usepackage{pdflscape}

% =====================================================
% Headers & Footers
% =====================================================
\usepackage{fancyhdr}
\pagestyle{fancy}
\fancyhf{}
\fancyfoot[C]{\thepage}
\begin{document}

\begin{center}

{\Large\textbf{ICS1512 Machine Learning Algorithms Laboratory}}

\vspace{1cm}
\begin{adjustbox}{width=\textwidth}
\begin{tabular}{ll}
\textbf{Experiment No.} & 1\\[0.3cm]

\textbf{Experiment Title} &
Working with Python Packages -- NumPy, SciPy, Scikit-Learn, Matplotlib\\[0.3cm]

\textbf{Student Name} & Sathiya Priyan\\[0.3cm]

\textbf{Register Number} & 3122247001058\\[0.3cm]

\textbf{Faculty} & Dr.\ Poreddy Ajay Kumar Reddy\\[0.3cm]

\textbf{Submission Date} & 30th June 2026

\end{tabular}
\end{adjustbox}
\end{center}

\vspace{1cm}

%%%%%%%%%%%%%%%%%%%%%%%%%%%%%%%%%%%%%%%%%%%%%%%%%%%%%%%
\section{Objective}

To explore the core functions and methods offered by five foundational Python
libraries used in Machine Learning:

\begin{itemize}
    \item NumPy
    \item Pandas
    \item SciPy
    \item Scikit-Learn
    \item Matplotlib
\end{itemize}

The experiment applies these libraries across a complete end-to-end machine
learning workflow.

Five heterogeneous public datasets were downloaded from Kaggle:

\begin{enumerate}
    \item Iris Dataset
    \item Pima Diabetes Dataset
    \item Loan Amount Prediction Dataset
    \item Email Spam Classification Dataset
    \item Handwritten Character Recognition Dataset
\end{enumerate}

For each dataset,

\begin{itemize}
    \item the appropriate machine learning task
          (classification or regression),
    \item a suitable feature-selection strategy,
    \item and an appropriate learning algorithm
\end{itemize}

were identified and empirically validated by benchmarking multiple models.

%%%%%%%%%%%%%%%%%%%%%%%%%%%%%%%%%%%%%%%%%%%%%%%%%%%%%%%
\section{Problem Statement}

Given five raw and unprocessed datasets that differ substantially in

\begin{itemize}
    \item size,
    \item dimensionality,
    \item and modality,
\end{itemize}

ranging from

\begin{itemize}
    \item a clean 4-feature numerical Iris dataset,
    \item a 1024-dimensional flattened handwritten image dataset,
    \item and a 3000-dimensional bag-of-words Email Spam dataset,
\end{itemize}

the objective is to perform the following tasks:

\begin{enumerate}[label=\textbf{(\roman*)}]
\item Load and explore each dataset using summary statistics and
visualizations.

\item Identify the correct machine learning task for each dataset:

\begin{itemize}
    \item Supervised Classification
    \item Supervised Regression
\end{itemize}

\item Perform suitable preprocessing:

\begin{itemize}
    \item Missing-value imputation
    \item Categorical encoding
    \item Feature scaling
    \item ANOVA-based feature selection for high-dimensional datasets
\end{itemize}

\item Split each dataset into training and testing partitions.

\item Train representative Scikit-Learn algorithms.

\item Evaluate and compare their performance using task-appropriate
performance metrics.
\end{enumerate}

The expected output of the experiment is

\begin{itemize}
    \item a task-identification table for all five datasets,
    \item together with quantitative performance benchmarks for every
    algorithm evaluated.
\end{itemize}

%%%%%%%%%%%%%%%%%%%%%%%%%%%%%%%%%%%%%%%%%%%%%%%%%%%%%%%%%%%%
\section{ML Task Identification}
%%%%%%%%%%%%%%%%%%%%%%%%%%%%%%%%%%%%%%%%%%%%%%%%%%%%%%%%%%%%

The appropriate machine learning task, feature-selection strategy, and the
empirically best-performing algorithm for each dataset are summarized in
Table~\ref{tab:ml_tasks}.

\begin{table}[H]
\centering
\caption{Machine Learning Task Identification}
\label{tab:ml_tasks}

\renewcommand{\arraystretch}{1.4}

\begin{adjustbox}{width=\textwidth}
\begin{tabular}{>{\raggedright\arraybackslash}p{3cm}
                >{\raggedright\arraybackslash}p{4cm}
                >{\raggedright\arraybackslash}p{4.2cm}
                >{\raggedright\arraybackslash}p{4.5cm}}
\toprule

\textbf{Dataset} &
\textbf{Type of ML Task} &
\textbf{Feature Selection Technique} &
\textbf{Suitable ML Algorithm (Empirically Best)} \\

\midrule

Iris Dataset &
Supervised, multi-class classification (3 classes) &
Not required -- only four low-noise, highly discriminative features &
Support Vector Machine (Linear Kernel), achieving 100\% test accuracy. \\

\midrule

Loan Amount Prediction &
Supervised regression (continuous target) &
Not applied -- only five numeric predictors remain after preprocessing; all were retained &
Gradient Boosting Regressor with the best performance,
$R^2 = 0.214$. \\

\midrule

Predicting Diabetes &
Supervised binary classification &
Not applied -- all eight clinical features are domain relevant &
Gradient Boosting / Random Forest,
approximately $75\%$ accuracy. \\

\midrule

Email Spam Classification &
Supervised binary classification with high-dimensional bag-of-words features &
SelectKBest using the ANOVA F-test, selecting the top $300$ features out of $3000$ &
Random Forest with
$96.04\%$ accuracy, ROC-AUC$=0.993$ \\

\midrule

Handwritten Character Recognition
(62 classes) &
Supervised multi-class classification using high-dimensional pixel features &
SelectKBest using the ANOVA F-test, selecting the top $100$ features out of $1024$ &
Random Forest with $35.19\%$ accuracy (best among the classical machine learning models tested). A Convolutional Neural Network (CNN) is recommended for production use. \\

\bottomrule
\end{tabular}
\end{adjustbox}
\end{table}

%%%%%%%%%%%%%%%%%%%%%%%%%%%%%%%%%%%%%%%%%%%%%%%%%%%%%%%%%%%%
\section{Dataset Description}
%%%%%%%%%%%%%%%%%%%%%%%%%%%%%%%%%%%%%%%%%%%%%%%%%%%%%%%%%%%%

The datasets used throughout this experiment differ significantly in terms of
size, dimensionality, feature types, and learning objectives.
Table~\ref{tab:dataset_description} summarizes the important
characteristics of every dataset.

\begin{table}[H]
\centering
\caption{Dataset Description}
\label{tab:dataset_description}

\renewcommand{\arraystretch}{1.35}

\begin{adjustbox}{width=\textwidth}
\begin{tabular}{>{\raggedright\arraybackslash}p{2.6cm}
                >{\raggedright\arraybackslash}p{3.8cm}
                >{\centering\arraybackslash}p{1.1cm}
                >{\centering\arraybackslash}p{1.2cm}
                >{\raggedright\arraybackslash}p{2.4cm}
                >{\raggedright\arraybackslash}p{2.4cm}
                >{\centering\arraybackslash}p{2cm}}

\toprule

\textbf{Dataset} &
\textbf{Source} &
\textbf{Samples} &
\textbf{Features} &
\textbf{Classes / Target} &
\textbf{Missing Values} &
\textbf{Train/Test Split} \\

\midrule

Iris &
Kaggle -- Iris Species Dataset &
150 &
4 Numeric &
3 Species &
None &
80\% / 20\% (Stratified) \\

\midrule

Diabetes &
Kaggle -- Pima Indians Diabetes Database &
768 &
8 Numeric &
Binary Outcome &
None &
80\% / 20\% (Stratified) \\

\midrule

Loan Amount Prediction &
Kaggle Loan Prediction Dataset
(\texttt{train\_u6lujuX\_CVtuZ9i.csv}) &
592 &
11 Mixed &
Continuous LoanAmount
(in thousands) &
Yes; Maximum missing percentage: $8.28\%$ in \texttt{Credit\_History} &
80\% / 20\% \\

\midrule

Email Spam &
Kaggle Email Spam Classification Dataset &
5172 &
3000 Word-count Features &
Binary
(Spam/Ham) &
None &
80\% / 20\% (Stratified) \\

\midrule

Handwritten Characters &
Kaggle Handwritten Character Recognition Dataset &
3410 &
1024
($32\times32$ grayscale pixels, flattened) &
62 Classes

(0--9, A--Z, a--z) &
None &
80\% / 20\% (Stratified) \\

\bottomrule

\end{tabular}
\end{adjustbox}

\end{table}

\vspace{0.3cm}

\noindent\textbf{Note:}

The original Loan Amount dataset contained 1614 records.
Rows having missing values in the target variable
(\texttt{LoanAmount}) were removed, leaving a final dataset of
592 samples for model training and evaluation.
\newpage
%%%%%%%%%%%%%%%%%%%%%%%%%%%%%%%%%%%%%%%%%%%%%%%%%%%%%%%%%%%%%%%%%%%%%%
\section{Exploratory Data Analysis}
%%%%%%%%%%%%%%%%%%%%%%%%%%%%%%%%%%%%%%%%%%%%%%%%%%%%%%%%%%%%%%%%%%%%%%

A single reusable \texttt{perform\_eda()} routine was developed to analyze every
dataset in a consistent manner. The routine computes:

\begin{itemize}
    \item Dataset shape
    \item Data types
    \item Summary statistics
    \item Missing-value counts
    \item Outlier counts using both:
    \begin{itemize}
        \item Interquartile Range (IQR)
        \item Z-score
    \end{itemize}
    \item Distribution plots
    \item Boxplots
    \item Correlation heatmaps
    \item Target-class distribution plots
\end{itemize}

For the two high-dimensional datasets

\begin{itemize}
    \item Email Spam Dataset (3000 features)
    \item Handwritten Character Dataset (1024 features)
\end{itemize}

the routine automatically skipped feature-wise histograms, boxplots, and
correlation heatmaps whenever the number of features exceeded 30, since
thousands of individual plots would be computationally expensive and difficult
to interpret. Only the target-class distribution was visualized for these
datasets.

%%%%%%%%%%%%%%%%%%%%%%%%%%%%%%%%%%%%%%%%%%%%%%%%%%%%%%%%%%%%%%%%%%%%%%
\subsection{Iris Dataset}
%%%%%%%%%%%%%%%%%%%%%%%%%%%%%%%%%%%%%%%%%%%%%%%%%%%%%%%%%%%%%%%%%%%%%%

\begin{figure}[htbp]
    \centering
    \includegraphics[width=\textwidth]{figures/iris_histograms.png}
    \caption{Feature-wise histograms of the Iris dataset (one-column layout, all four numeric features).}
    \label{fig:iris_histograms}
\end{figure}

\subsubsection*{Inference}

All four Iris features are either unimodal or bimodal with no missing values.
Among them,

\begin{itemize}
    \item Petal Length
    \item Petal Width
\end{itemize}

show a visibly bimodal (or trimodal) distribution, indicating that these
features separate the three Iris species much more effectively than the sepal
measurements, whose distributions overlap considerably between
\textit{Iris-versicolor} and \textit{Iris-virginica}.

This observation agrees with the well-known characteristic of the Iris dataset,
where petal measurements are historically the strongest discriminative features.

Sepal Width appears to be the most symmetric and approximately Gaussian feature
among the four.

%%%%%%%%%%%%%%%%%%%%%%%%%%%%%%%%%%%%%%%%%%%%%%%%%%%%%%%%%%%%%%%%%%%%%%

\begin{figure}[H]
\centering

\begin{minipage}{0.48\textwidth}
    \centering
    \includegraphics[width=\linewidth]{figures/iris_boxplot.png}
    \caption{Boxplots of the four Iris features.}
    \label{fig:iris_boxplot}
\end{minipage}
\hfill
\begin{minipage}{0.48\textwidth}
    \centering
    \includegraphics[width=\linewidth]{figures/iris_heatmap.png}
    \caption{Correlation heatmap of the Iris features.}
    \label{fig:iris_heatmap}
\end{minipage}

\end{figure}

\subsubsection*{Inference}

The boxplots reveal only a few mild outliers in
\texttt{SepalWidthCm}:

\begin{itemize}
    \item 4 outliers detected using the IQR rule
    \item 1 outlier detected using the stricter
          $|z|>3$ criterion
\end{itemize}

No meaningful outliers are present in the remaining three features,
confirming that the Iris dataset is essentially clean.

The correlation heatmap shows that

\begin{itemize}
    \item Petal Length and Petal Width are almost perfectly correlated
          ($r \approx 0.96$).

    \item Petal Length is also strongly correlated with Sepal Length
          ($r \approx 0.87$).

    \item Petal Width is strongly correlated with Sepal Length
          ($r \approx 0.82$).

    \item Sepal Width exhibits weak and slightly negative correlations
          with the remaining features.
\end{itemize}

The strong redundancy between Petal Length and Petal Width explains why a
linear decision boundary is capable of separating the Iris classes almost
perfectly.

%%%%%%%%%%%%%%%%%%%%%%%%%%%%%%%%%%%%%%%%%%%%%%%%%%%%%%%%%%%%%%%%%%%%%%

\begin{figure}[H]
    \centering
    \includegraphics[width=0.6\textwidth]{figures/iris_target_distribution.png}
    \caption{Class distribution of the target variable \texttt{Species}.}
    \label{fig:iris_target}
\end{figure}

\subsubsection*{Inference}

The target variable is perfectly balanced, containing exactly

\[
50
\]

samples from each class:

\begin{itemize}
    \item Iris-setosa
    \item Iris-versicolor
    \item Iris-virginica
\end{itemize}

Since every class contains the same number of observations, no class-imbalance
correction techniques such as

\begin{itemize}
    \item Oversampling
    \item Undersampling
    \item Class weighting
\end{itemize}

are required.

Consequently, classification accuracy is an appropriate evaluation metric for
this dataset because each class contributes equally to the overall performance.

%%%%%%%%%%%%%%%%%%%%%%%%%%%%%%%%%%%%%%%%%%%%%%%%%%%%%%%%%%%%%%%%%%%%%%
\subsection{Diabetes Dataset}
%%%%%%%%%%%%%%%%%%%%%%%%%%%%%%%%%%%%%%%%%%%%%%%%%%%%%%%%%%%%%%%%%%%%%%

\begin{figure}[htbp]
    \centering
    \includegraphics[width=\textwidth]{figures/diabetes_histograms.png}
    \caption{Feature-wise histograms of the Pima Diabetes dataset (one-column layout, all eight numeric features).}
    \label{fig:diabetes_histograms}
\end{figure}

\subsubsection*{Inference}

The feature distributions reveal several important characteristics of the
dataset.

\begin{itemize}

\item \textbf{Insulin} and
\textbf{DiabetesPedigreeFunction} exhibit heavily right-skewed
distributions with long tails.

\item \textbf{Glucose},
\textbf{BloodPressure}, and
\textbf{BMI} appear approximately normally distributed; however,
each contains a noticeable spike at the value zero.

\item Since a living patient cannot realistically have

\begin{itemize}
    \item zero blood pressure,
    \item zero body mass index,
    \item or zero glucose level,
\end{itemize}

these zero values most likely represent disguised missing values rather
than valid physiological measurements.

\item The \textbf{Age} variable is also right-skewed,
indicating that the dataset contains more younger patients than older
patients.

\item These zero-inflated distributions contribute directly to the
large number of outliers detected using both the IQR and Z-score methods.

\end{itemize}

%%%%%%%%%%%%%%%%%%%%%%%%%%%%%%%%%%%%%%%%%%%%%%%%%%%%%%%%%%%%%%%%%%%%%%

\begin{figure}[H]
\centering

\begin{minipage}{0.48\textwidth}
    \centering
    \includegraphics[width=\linewidth]{figures/diabetes_boxplot.png}
    \caption{Boxplots of the eight Diabetes features.}
    \label{fig:diabetes_boxplot}
\end{minipage}
\hfill
\begin{minipage}{0.48\textwidth}
    \centering
    \includegraphics[width=\linewidth]{figures/diabetes_heatmap.png}
    \caption{Correlation heatmap of the Diabetes features.}
    \label{fig:diabetes_heatmap}
\end{minipage}

\end{figure}

\subsubsection*{Inference}

The boxplots indicate that several variables contain a considerable number of
outliers.

\begin{itemize}

\item \textbf{BloodPressure} has the largest number of detected
outliers:

\begin{itemize}
    \item 45 outliers using the IQR rule.
    \item 35 outliers using the Z-score criterion.
\end{itemize}

These are primarily caused by the disguised zero values identified
earlier.

\item \textbf{Insulin} contains 34 IQR outliers because of its highly
right-skewed distribution.

\end{itemize}

The correlation heatmap provides additional insights into the relationships
among the predictor variables.

\begin{itemize}

\item \textbf{Glucose} exhibits the strongest positive correlation with
the target variable \texttt{Outcome}, which agrees with medical
knowledge since elevated blood glucose is the primary indicator of
diabetes.

\item \textbf{Age} and \textbf{Pregnancies} are moderately positively
correlated, which is biologically expected.

\item No pair of predictor variables exhibits severe multicollinearity.

\end{itemize}

Therefore, multicollinearity is not expected to significantly affect the
performance of the linear models evaluated in this experiment.

%%%%%%%%%%%%%%%%%%%%%%%%%%%%%%%%%%%%%%%%%%%%%%%%%%%%%%%%%%%%%%%%%%%%%%

\begin{figure}[H]
    \centering
    \includegraphics[width=0.65\textwidth]{figures/diabetes_target_distribution.png}
    \caption{Class distribution of the target variable \texttt{Outcome}.}
    \label{fig:diabetes_target}
\end{figure}

\subsubsection*{Inference}

The target variable is moderately imbalanced.

Approximately

\[
65\%
\]

of the observations correspond to

\[
\texttt{Outcome}=0
\]

(non-diabetic patients), while approximately

\[
35\%
\]

belong to

\[
\texttt{Outcome}=1
\]

(diabetic patients).

This observation is consistent with the mean value of the target variable,

\[
\overline{\texttt{Outcome}} = 0.349,
\]

reported in the summary statistics.

Because of this class imbalance, model performance was evaluated using
multiple classification metrics instead of relying solely on overall
accuracy.

The reported evaluation metrics include

\begin{itemize}

\item Accuracy

\item Precision

\item Recall

\item Weighted F1-score

\end{itemize}

These metrics provide a more reliable assessment of classifier performance
under imbalanced class distributions.

%%%%%%%%%%%%%%%%%%%%%%%%%%%%%%%%%%%%%%%%%%%%%%%%%%%%%%%%%%%%%%%%%%%%%%
\subsection{Loan Amount Prediction Dataset}
%%%%%%%%%%%%%%%%%%%%%%%%%%%%%%%%%%%%%%%%%%%%%%%%%%%%%%%%%%%%%%%%%%%%%%

\begin{figure}[H]
    \centering
    \includegraphics[width=\textwidth]{figures/loan_missing_heatmap.png}
    \caption{Missing-value heatmap of the Loan Amount dataset (yellow indicates missing values).}
    \label{fig:loan_missing_heatmap}
\end{figure}

\subsubsection*{Inference}

Unlike the other four datasets, the Loan Amount Prediction dataset contains
genuine missing values distributed across multiple features.

The highest percentage of missing values is observed in:

\begin{itemize}
    \item \texttt{Credit\_History} -- $8.28\%$
    \item \texttt{Self\_Employed} -- $5.24\%$
    \item \texttt{Loan\_Amount\_Term} -- $2.36\%$
    \item \texttt{Gender} -- approximately $2.2\%$
    \item \texttt{Dependents} -- approximately $2.2\%$
    \item \texttt{Married} -- $0.34\%$
\end{itemize}

The missing values are scattered throughout the dataset rather than being
concentrated within a particular group of rows or columns.

This pattern suggests that the data are most likely
\textbf{Missing At Random (MAR)}, where applicants omitted optional
information rather than the data being lost due to systematic collection
errors.

%%%%%%%%%%%%%%%%%%%%%%%%%%%%%%%%%%%%%%%%%%%%%%%%%%%%%%%%%%%%%%%%%%%%%%

\begin{figure}[H]
    \centering
    \includegraphics[width=\textwidth]{figures/loan_histograms.png}
    \caption{Feature-wise histograms of the Loan Amount dataset (numeric features).}
    \label{fig:loan_histograms}
\end{figure}

\subsubsection*{Inference}

The feature distributions exhibit several important characteristics.

\begin{itemize}

\item \textbf{ApplicantIncome} and
\textbf{CoapplicantIncome} are highly right-skewed.

Only a small number of applicants possess extremely high incomes,
creating long tails in both distributions.

These high-income observations correspond directly to the

\begin{itemize}
    \item 49 IQR outliers for ApplicantIncome, and
    \item 18 IQR outliers for CoapplicantIncome.
\end{itemize}

\item \textbf{Loan\_Amount\_Term} is dominated by a single value of

\[
360 \text{ months}
\]

which represents the standard 30-year loan tenure, while only a few
applications correspond to shorter repayment periods.

\item \textbf{Credit\_History} behaves as a binary variable rather than a
continuous numerical feature.

Consequently, the IQR method incorrectly identifies approximately
85 observations as outliers.

This is a well-known limitation of applying outlier detection techniques
designed for continuous variables to near-binary data.

\end{itemize}

%%%%%%%%%%%%%%%%%%%%%%%%%%%%%%%%%%%%%%%%%%%%%%%%%%%%%%%%%%%%%%%%%%%%%%

\begin{figure}[H]
\centering

\begin{minipage}{0.48\textwidth}
    \centering
    \includegraphics[width=\linewidth]{figures/loan_boxplot.png}
    \caption{Boxplots of the Loan Amount numeric features.}
    \label{fig:loan_boxplot}
\end{minipage}
\hfill
\begin{minipage}{0.48\textwidth}
    \centering
    \includegraphics[width=\linewidth]{figures/loan_heatmap.png}
    \caption{Correlation heatmap of the Loan Amount numeric features.}
    \label{fig:loan_heatmap}
\end{minipage}

\end{figure}

\subsubsection*{Inference}

The boxplots clearly illustrate that income-related outliers dominate the
scale of the remaining numerical variables.

Because all variables are plotted on the same axis,

\begin{itemize}
    \item ApplicantIncome and CoapplicantIncome exhibit extremely large
          ranges.

    \item Credit\_History appears as a thin bar close to the value

\[
1,
\]

since nearly

\[
84\%
\]

of applicants possess a positive credit history.

\end{itemize}

The correlation heatmap indicates only weak pairwise relationships among
the remaining numerical predictors.

No feature exhibits a strong linear relationship either

\begin{itemize}
    \item with another predictor, or

    \item with the regression target,
          \texttt{LoanAmount}.
\end{itemize}

This weak correlation foreshadows the relatively poor regression
performance discussed later.

%%%%%%%%%%%%%%%%%%%%%%%%%%%%%%%%%%%%%%%%%%%%%%%%%%%%%%%%%%%%%%%%%%%%%%

\begin{figure}[H]
    \centering
    \includegraphics[width=0.75\textwidth]{figures/loan_target_distribution.png}
    \caption{Distribution of the regression target \texttt{LoanAmount} (in thousands).}
    \label{fig:loan_target}
\end{figure}

\subsubsection*{Inference}

The target variable \texttt{LoanAmount} exhibits a pronounced
right-skewed distribution.

Most approved loans are concentrated between approximately

\[
100
\]

and

\[
170
\]

(thousand), while relatively few loans extend to values as high as

\[
700.
\]

This skewness causes ordinary linear regression models to be
disproportionately influenced by a small number of large-loan observations.

Consequently, tree-based ensemble methods such as

\begin{itemize}
    \item Random Forest Regressor
    \item Gradient Boosting Regressor
\end{itemize}

outperform Linear Regression, Ridge Regression, and Lasso Regression in
the final experimental comparison.

%%%%%%%%%%%%%%%%%%%%%%%%%%%%%%%%%%%%%%%%%%%%%%%%%%%%%%%%%%%%%%%%%%%%%%
\subsection{Email Spam Dataset}
\label{subsec:email_spam}
%%%%%%%%%%%%%%%%%%%%%%%%%%%%%%%%%%%%%%%%%%%%%%%%%%%%%%%%%%%%%%%%%%%%%%

\subsubsection*{High-Dimensional Dataset -- Detailed Feature Plots Skipped}

The Email Spam dataset contains approximately \textbf{3000 word-count
features}, making it a high-dimensional dataset. To maintain readability and
avoid generating an excessive number of plots, the
\texttt{perform\_eda()} routine automatically skipped feature-wise
histograms, boxplots, and correlation heatmaps.

Instead, only the target-class distribution was visualized, while summary
statistics and outlier information were printed to the console.

%%%%%%%%%%%%%%%%%%%%%%%%%%%%%%%%%%%%%%%%%%%%%%%%%%%%%%%%%%%%%%%%%%%%%%

\begin{figure}[H]
    \centering
    \includegraphics[width=0.65\textwidth]{figures/email_target_distribution.png}
    \caption{Class distribution of the target variable \texttt{Prediction}
    (0 = Ham, 1 = Spam).}
    \label{fig:email_target_distribution}
\end{figure}

%%%%%%%%%%%%%%%%%%%%%%%%%%%%%%%%%%%%%%%%%%%%%%%%%%%%%%%%%%%%%%%%%%%%%%

\subsubsection*{Inference}

Since the dataset contains

\[
3000
\]

word-count features, generating individual feature plots would require

\begin{itemize}
    \item 3000 histograms,
    \item 3000 boxplots,
    \item and a
    \[
    3000 \times 3000
    \]
    correlation matrix.
\end{itemize}

Such visualizations would be computationally expensive and practically
impossible to interpret.

Therefore, the exploratory data analysis routine automatically suppressed
these plots and generated only the target-class distribution.

The class distribution indicates a moderate imbalance between the two
classes.

Approximately

\[
71\%
\]

of the emails belong to the

\[
\texttt{Ham}
\]

class, whereas approximately

\[
29\%
\]

belong to the

\[
\texttt{Spam}
\]

class.

This observation is consistent with the reported mean value of the
\texttt{Prediction} column,

\[
\overline{\texttt{Prediction}} = 0.290.
\]

Because of this class imbalance, model evaluation relied on multiple
performance metrics rather than overall accuracy alone.

The reported metrics include

\begin{itemize}
    \item Accuracy
    \item Precision
    \item Recall
    \item Weighted F1-score
\end{itemize}

to provide a more reliable assessment of classifier performance.

The exploratory statistics also revealed that the most frequently occurring
word-count features, including

\begin{itemize}
    \item \texttt{the}
    \item \texttt{to}
    \item \texttt{ect}
    \item \texttt{and}
    \item \texttt{for}
\end{itemize}

have extremely large occurrence counts and correspondingly high IQR outlier
counts, ranging approximately from

\[
400
\]

to

\[
665
\]

observations.

This behavior is expected because these words are among the most common
English stop words and naturally appear in a large proportion of emails.

Such high-frequency features contribute relatively little toward
distinguishing spam from legitimate emails.

Consequently, applying feature selection using

\[
\texttt{SelectKBest}(\texttt{f\_classif})
\]

to retain only the most discriminative features is an effective strategy for
reducing dimensionality while preserving classification performance.

%%%%%%%%%%%%%%%%%%%%%%%%%%%%%%%%%%%%%%%%%%%%%%%%%%%%%%%%%%%%%%%%%%%%%%
\subsection{Handwritten Character Recognition Dataset}
\label{subsec:handwritten_characters}
%%%%%%%%%%%%%%%%%%%%%%%%%%%%%%%%%%%%%%%%%%%%%%%%%%%%%%%%%%%%%%%%%%%%%%

\subsubsection*{High-Dimensional Dataset -- Detailed Feature Plots Skipped}

The Handwritten Character Recognition dataset consists of

\[
1024
\]

pixel features obtained by flattening each

\[
32 \times 32
\]

grayscale image into a one-dimensional feature vector.

Since the dataset contains a very large number of features, the
\texttt{perform\_eda()} routine automatically skipped the generation of

\begin{itemize}
    \item Feature-wise histograms,
    \item Boxplots,
    \item Correlation heatmaps.
\end{itemize}

Instead, only the target-class distribution was visualized, while summary
statistics and outlier information were reported in the console.

%%%%%%%%%%%%%%%%%%%%%%%%%%%%%%%%%%%%%%%%%%%%%%%%%%%%%%%%%%%%%%%%%%%%%%

\begin{figure}[H]
    \centering
    \includegraphics[width=\textwidth]{figures/handwritten_target_distribution.png}
    \caption{Class distribution across the 62 handwritten character labels (0--9, A--Z, a--z).}
    \label{fig:handwritten_target_distribution}
\end{figure}

%%%%%%%%%%%%%%%%%%%%%%%%%%%%%%%%%%%%%%%%%%%%%%%%%%%%%%%%%%%%%%%%%%%%%%

\subsubsection*{Inference}

The dataset contains

\[
62
\]

distinct output classes representing

\begin{itemize}
    \item Digits (0--9),
    \item Uppercase letters (A--Z),
    \item Lowercase letters (a--z).
\end{itemize}

Since each image is represented using

\[
1024
\]

flattened pixel values, detailed feature-wise visualizations were omitted
for the same readability and computational reasons discussed for the Email
Spam dataset.

The target-class distribution is relatively uniform across all
62 character classes.

Each class contains approximately

\[
50\text{--}60
\]

samples, indicating that class imbalance is not a significant concern for
this dataset.

Instead, the primary challenge lies in the extremely high dimensionality of
the feature space.

More importantly, flattening each

\[
32 \times 32
\]

image into a one-dimensional vector destroys the inherent two-dimensional
spatial relationships between neighboring pixels.

As a result, traditional machine learning algorithms operate on independent
pixel values rather than exploiting local spatial structures.

Consequently, classical classifiers such as

\begin{itemize}
    \item Logistic Regression,
    \item Support Vector Machines,
    \item Decision Trees,
    \item Random Forests,
    \item Gradient Boosting,
\end{itemize}

are unable to fully capture the visual patterns present in handwritten
characters.

This limitation is discussed further in the Discussion section, where the
relatively lower classification accuracy is attributed primarily to the
loss of spatial information rather than deficiencies in the learning
algorithms themselves.

%%%%%%%%%%%%%%%%%%%%%%%%%%%%%%%%%%%%%%%%%%%%%%%%%%%%%%%%%%%%%%%%%%%%%%
\section{Data Preprocessing}
%%%%%%%%%%%%%%%%%%%%%%%%%%%%%%%%%%%%%%%%%%%%%%%%%%%%%%%%%%%%%%%%%%%%%%

A single reusable preprocessing pipeline,
\texttt{run\_ml\_pipeline()}, was applied uniformly to all five datasets to
ensure consistency throughout the experiment.

The preprocessing steps are summarized below.

\begin{enumerate}

\item \textbf{Missing Value Handling}

Numeric attributes were imputed using the column median through

\[
\texttt{SimpleImputer(strategy='median')}.
\]

Categorical attributes were imputed using the most frequently occurring
value through

\[
\texttt{SimpleImputer(strategy='most\_frequent')}.
\]

Among the five datasets, only the Loan Amount Prediction dataset required
missing-value imputation, since the remaining four datasets contained no
missing values.

\vspace{0.3cm}

\item \textbf{Categorical Encoding}

Every categorical feature was transformed into numerical form using

\[
\texttt{LabelEncoder}.
\]

Similarly, non-numeric target variables such as

\begin{itemize}
    \item Species,
    \item Prediction,
    \item Outcome,
    \item Label
\end{itemize}

were encoded before model training to ensure compatibility with
Scikit-Learn classifiers.

\vspace{0.3cm}

\item \textbf{Feature Scaling}

All numerical predictor variables were standardized using

\[
\texttt{StandardScaler},
\]

which transforms every feature according to

\[
z=\frac{x-\mu}{\sigma},
\]

where

\begin{itemize}
    \item $x$ is the original feature value,
    \item $\mu$ is the feature mean,
    \item $\sigma$ is the standard deviation.
\end{itemize}

This produces standardized features having

\[
\mu=0,\qquad
\sigma=1.
\]

Feature scaling was performed for every dataset before model training.

\vspace{0.3cm}

\item \textbf{Train--Test Split}

Every dataset was divided into

\[
80\% \text{ training}
\]

and

\[
20\% \text{ testing}
\]

using

\[
\texttt{test\_size}=0.2,
\qquad
\texttt{random\_state}=42.
\]

For all classification datasets,

\[
\texttt{stratify}=y
\]

was enabled to preserve the original class distribution in both the
training and testing sets.

The Loan Amount Prediction dataset is a regression problem; therefore,
stratified sampling was not applied.

\vspace{0.3cm}

\item \textbf{Feature Engineering / Feature Selection}

Feature selection was applied only to high-dimensional datasets.

\begin{itemize}

\item \textbf{Email Spam Dataset}

The original

\[
3000
\]

word-count features were reduced to the

\[
300
\]

most informative features using

\[
\texttt{SelectKBest}(f_{\text{classif}}).
\]

\vspace{0.2cm}

\item \textbf{Handwritten Character Dataset}

The original

\[
1024
\]

pixel features were reduced to the top

\[
100
\]

features using

\[
\texttt{SelectKBest}(f_{\text{classif}}).
\]

\vspace{0.2cm}

\item \textbf{Iris, Diabetes, and Loan Amount}

No feature selection was applied because these datasets already contain
relatively few predictors, all of which were retained for model training.

\end{itemize}

\end{enumerate}

%%%%%%%%%%%%%%%%%%%%%%%%%%%%%%%%%%%%%%%%%%%%%%%%%%%%%%%%%%%%%%%%%%%%%%
\section{Results}
%%%%%%%%%%%%%%%%%%%%%%%%%%%%%%%%%%%%%%%%%%%%%%%%%%%%%%%%%%%%%%%%%%%%%%

Five datasets and up to eight machine learning algorithms were evaluated
during the experiment.

Since the datasets differ considerably in their characteristics and machine
learning objectives, the experimental results are presented separately for
each dataset instead of combining them into a single comparison table.

All numerical values reported in the following sections are taken directly
from the captured output of the Jupyter Notebook.

%%%%%%%%%%%%%%%%%%%%%%%%%%%%%%%%%%%%%%%%%%%%%%%%%%%%%%%%%%%%%%%%%%%%%%
\subsection{Iris Dataset (Classification)}
%%%%%%%%%%%%%%%%%%%%%%%%%%%%%%%%%%%%%%%%%%%%%%%%%%%%%%%%%%%%%%%%%%%%%%

\begin{table}[H]
\centering
\caption{Performance Comparison on the Iris Dataset}
\label{tab:iris_results}

\renewcommand{\arraystretch}{1.2}

\begin{adjustbox}{width=\textwidth}
\begin{tabular}{|l|c|c|c|c|c|}

\toprule

\textbf{Model}
&
\textbf{Accuracy}
&
\textbf{Precision}
&
\textbf{Recall}
&
\textbf{F1-Score}
&
\textbf{ROC-AUC}

\\

\midrule

Logistic Regression &
0.9333 &
0.9333 &
0.9333 &
0.9333 &
0.9967 \\

K-Nearest Neighbors &
0.9333 &
0.9444 &
0.9333 &
0.9327 &
0.9933 \\

SVM (RBF Kernel) &
0.9667 &
0.9697 &
0.9667 &
0.9666 &
0.9967 \\

SVM (Linear Kernel) &
\textbf{1.0000} &
\textbf{1.0000} &
\textbf{1.0000} &
\textbf{1.0000} &
\textbf{1.0000} \\

Decision Tree &
0.9000 &
0.9024 &
0.9000 &
0.8997 &
0.9250 \\

Random Forest &
0.9000 &
0.9024 &
0.9000 &
0.8997 &
0.9933 \\

Gradient Boosting &
0.9000 &
0.9024 &
0.9000 &
0.8997 &
0.9933 \\

Gaussian Naive Bayes &
0.9667 &
0.9697 &
0.9667 &
0.9666 &
0.9900 \\

\bottomrule

\end{tabular}
\end{adjustbox}

\end{table}

%%%%%%%%%%%%%%%%%%%%%%%%%%%%%%%%%%%%%%%%%%%%%%%%%%%%%%%%%%%%%%%%%%%%%%
\subsection{Diabetes Dataset (Classification)}
%%%%%%%%%%%%%%%%%%%%%%%%%%%%%%%%%%%%%%%%%%%%%%%%%%%%%%%%%%%%%%%%%%%%%%

\begin{table}[H]
\centering
\caption{Performance Comparison on the Diabetes Dataset}
\label{tab:diabetes_results}

\renewcommand{\arraystretch}{1.2}

\begin{adjustbox}{width=\textwidth}
\begin{tabular}{|l|c|c|c|c|c|}

\toprule

\textbf{Model}
&
\textbf{Accuracy}
&
\textbf{Precision}
&
\textbf{Recall}
&
\textbf{F1-Score}
&
\textbf{ROC-AUC}

\\

\midrule

Logistic Regression &
0.7143 &
0.7065 &
0.7143 &
0.7084 &
0.8230 \\

K-Nearest Neighbors &
0.7078 &
0.7023 &
0.7078 &
0.7043 &
0.7458 \\

SVM (RBF Kernel) &
0.7468 &
0.7438 &
0.7468 &
0.7450 &
0.7933 \\

SVM (Linear Kernel) &
0.7208 &
0.7126 &
0.7208 &
0.7141 &
0.8275 \\

Decision Tree &
0.7208 &
0.7108 &
0.7208 &
0.7102 &
0.6657 \\

Random Forest &
0.7532 &
0.7497 &
0.7532 &
0.7509 &
0.8145 \\

Gradient Boosting &
\textbf{0.7532} &
0.7483 &
0.7532 &
0.7496 &
\textbf{0.8394} \\

Gaussian Naive Bayes &
0.7078 &
0.7179 &
0.7078 &
0.7114 &
0.7728 \\

\bottomrule

\end{tabular}
\end{adjustbox}

\end{table}

%%%%%%%%%%%%%%%%%%%%%%%%%%%%%%%%%%%%%%%%%%%%%%%%%%%%%%%%%%%%%%%%%%%%%%
\subsection{Loan Amount Prediction (Regression)}
%%%%%%%%%%%%%%%%%%%%%%%%%%%%%%%%%%%%%%%%%%%%%%%%%%%%%%%%%%%%%%%%%%%%%%

\begin{table}[H]
\centering
\caption{Performance Comparison on the Loan Amount Prediction Dataset}
\label{tab:loan_results}

\renewcommand{\arraystretch}{1.2}

\begin{adjustbox}{width=\textwidth}
\begin{tabular}{lcccccc}

\toprule

\textbf{Model}
&
\textbf{MAE}
&
\textbf{MSE}
&
\textbf{RMSE}
&
$\mathbf{R^2}$
&
\textbf{MAPE (\%)}
&
\textbf{Adjusted $R^2$}

\\

\midrule

Linear Regression
&
36.6989
&
3324.82
&
57.6612
&
0.0981
&
35.4658
&
0.0054
\\

Ridge Regression
&
36.7015
&
3322.14
&
57.6380
&
0.0988
&
35.4790
&
0.0062
\\

Lasso Regression
&
36.1555
&
3273.65
&
57.2158
&
0.1120
&
35.3014
&
0.0207
\\

Decision Tree Regressor
&
48.9832
&
5382.09
&
73.3627
&
--0.4600
&
42.2095
&
--0.6101
\\

Random Forest Regressor
&
33.9438
&
3132.68
&
55.9703
&
0.1502
&
30.5887
&
0.0628
\\

Gradient Boosting Regressor
&
\textbf{33.4631}
&
\textbf{2897.25}
&
\textbf{53.8261}
&
\textbf{0.2141}
&
\textbf{29.7789}
&
\textbf{0.1333}
\\

Support Vector Regression (SVR)
&
39.5036
&
3517.44
&
59.3080
&
0.0458
&
36.2423
&
--0.0523
\\

\bottomrule

\end{tabular}
\end{adjustbox}

\end{table}

%%%%%%%%%%%%%%%%%%%%%%%%%%%%%%%%%%%%%%%%%%%%%%%%%%%%%%%%%%%%%%%%%%%%%%
\subsection{Email Spam Classification}
%%%%%%%%%%%%%%%%%%%%%%%%%%%%%%%%%%%%%%%%%%%%%%%%%%%%%%%%%%%%%%%%%%%%%%

\begin{table}[H]
\centering
\caption{Performance Comparison on the Email Spam Dataset}
\label{tab:spam_results}

\renewcommand{\arraystretch}{1.2}

\begin{adjustbox}{width=\textwidth}
\begin{tabular}{|l|c|c|c|c|c|}

\toprule

\textbf{Model}
&
\textbf{Accuracy}
&
\textbf{Precision}
&
\textbf{Recall}
&
\textbf{F1-Score}
&
\textbf{ROC-AUC}

\\

\midrule

Logistic Regression
&
0.9459
&
0.9470
&
0.9459
&
0.9463
&
0.9871
\\

K-Nearest Neighbors
&
0.9266
&
0.9261
&
0.9266
&
0.9263
&
0.9727
\\

SVM (RBF Kernel)
&
0.9353
&
0.9347
&
0.9353
&
0.9347
&
0.9864
\\

SVM (Linear Kernel)
&
0.9382
&
0.9407
&
0.9382
&
0.9388
&
0.9842
\\

Decision Tree
&
0.9246
&
0.9243
&
0.9246
&
0.9244
&
0.9055
\\

Random Forest
&
\textbf{0.9604}
&
\textbf{0.9603}
&
\textbf{0.9604}
&
\textbf{0.9603}
&
\textbf{0.9932}
\\

Gradient Boosting
&
0.9575
&
0.9577
&
0.9575
&
0.9576
&
0.9904
\\

Gaussian Naive Bayes
&
0.8705
&
0.8799
&
0.8705
&
0.8603
&
0.9747
\\

\bottomrule

\end{tabular}
\end{adjustbox}

\end{table}

%%%%%%%%%%%%%%%%%%%%%%%%%%%%%%%%%%%%%%%%%%%%%%%%%%%%%%%%%%%%%%%%%%%%%%
\subsection{Handwritten Character Recognition}
%%%%%%%%%%%%%%%%%%%%%%%%%%%%%%%%%%%%%%%%%%%%%%%%%%%%%%%%%%%%%%%%%%%%%%

\begin{table}[H]
\centering
\caption{Performance Comparison on the Handwritten Character Recognition Dataset}
\label{tab:handwritten_results}

\renewcommand{\arraystretch}{1.2}

\begin{adjustbox}{width=\textwidth}
\begin{tabular}{|l|c|c|c|c|c|}

\toprule

\textbf{Model}
&
\textbf{Accuracy}
&
\textbf{Precision}
&
\textbf{Recall}
&
\textbf{F1-Score}
&
\textbf{ROC-AUC}

\\

\midrule

Logistic Regression
&
0.2478
&
0.2656
&
0.2478
&
0.2468
&
0.8347
\\

K-Nearest Neighbors
&
0.2933
&
0.3283
&
0.2933
&
0.2892
&
0.7862
\\

SVM (RBF Kernel)
&
0.3460
&
0.3818
&
0.3460
&
0.3467
&
0.8987
\\

SVM (Linear Kernel)
&
0.2903
&
0.2966
&
0.2903
&
0.2853
&
0.8618
\\

Decision Tree
&
0.2023
&
0.2275
&
0.2023
&
0.2082
&
0.5959
\\

Random Forest
&
\textbf{0.3519}
&
0.3569
&
\textbf{0.3519}
&
0.3464
&
0.8981
\\

Gradient Boosting
&
0.2390
&
0.2503
&
0.2390
&
0.2359
&
0.8320
\\

Gaussian Naive Bayes
&
0.1188
&
0.2285
&
0.1188
&
0.1062
&
0.7527
\\

\bottomrule

\end{tabular}
\end{adjustbox}

\end{table}

%%%%%%%%%%%%%%%%%%%%%%%%%%%%%%%%%%%%%%%%%%%%%%%%%%%%%%%%%%%%%%%%%%%%%%
\section{Discussion}
%%%%%%%%%%%%%%%%%%%%%%%%%%%%%%%%%%%%%%%%%%%%%%%%%%%%%%%%%%%%%%%%%%%%%%

The experimental results demonstrate that model performance depends far
more strongly on dataset characteristics than on the choice of machine
learning algorithm alone.

The Iris dataset represents the easiest learning problem.

Its four clean and highly discriminative features allow even relatively
simple models to achieve excellent performance.

The linear-kernel Support Vector Machine achieved perfect

\[
100\%
\]

classification accuracy, while even the weakest models exceeded

\[
90\%.
\]

The Diabetes dataset proved considerably more challenging.

Most algorithms achieved approximately

\[
75\%
\]

accuracy, largely because of the disguised missing values represented as
zeros and the natural overlap between diabetic and non-diabetic patient
measurements.

The Loan Amount Prediction dataset produced the weakest regression
results.

The best-performing model,

Gradient Boosting Regressor,

achieved only

\[
R^2=0.214,
\]

indicating that the available predictor variables provide limited
information regarding the approved loan amount.

In practical banking applications, loan approval depends upon many
additional variables that are absent from this dataset.

The Email Spam dataset achieved excellent classification performance.

Random Forest obtained

\[
96.04\%
\]

accuracy together with a

\[
0.993
\]

ROC-AUC score.

The strong performance can be attributed to the highly discriminative
nature of word-frequency features for spam detection.

The Handwritten Character Recognition dataset was the most difficult
classification task.

Although Random Forest achieved the highest accuracy among the evaluated
models,

\[
35.19\%
\]

remains substantially below the performance expected from modern deep
learning methods.

The primary limitation arises from representing each

\[
32\times32
\]

image as a flattened vector, thereby discarding spatial information that
Convolutional Neural Networks naturally exploit.

Across all five datasets, tree-based ensemble methods such as Random
Forest and Gradient Boosting consistently ranked among the strongest
performers.

Conversely, Gaussian Naive Bayes generally performed poorly whenever its
conditional independence assumption was violated.

%%%%%%%%%%%%%%%%%%%%%%%%%%%%%%%%%%%%%%%%%%%%%%%%%%%%%%%%%%%%%%%%%%%%%%
\section{Conclusion}
%%%%%%%%%%%%%%%%%%%%%%%%%%%%%%%%%%%%%%%%%%%%%%%%%%%%%%%%%%%%%%%%%%%%%%

This experiment successfully explored the functionality of the Python
libraries NumPy, Pandas, SciPy, Scikit-Learn, and Matplotlib through a
single reusable machine learning pipeline.

Five structurally different public datasets were analyzed using a common
workflow consisting of

\begin{itemize}
    \item Exploratory Data Analysis,
    \item Data Preprocessing,
    \item Feature Selection,
    \item Model Training,
    \item Performance Evaluation.
\end{itemize}

The correct machine learning task was identified for every dataset, and
appropriate feature-selection techniques were applied to
high-dimensional datasets using the ANOVA F-test through
\texttt{SelectKBest}.

Eight classification algorithms and seven regression algorithms were
benchmarked using a common evaluation framework.

The results demonstrate that classical machine learning algorithms

\begin{itemize}
    \item perform exceptionally well on clean, low-dimensional datasets
          such as Iris,

    \item achieve strong performance on high-dimensional text data such
          as Email Spam,

    \item obtain moderate performance on noisy clinical datasets such as
          Diabetes,

    \item perform relatively poorly when important predictive variables
          are absent, as in Loan Amount Prediction,

    \item and are fundamentally limited when applied directly to
          flattened image pixels instead of spatial image
          representations.
\end{itemize}

Overall, the experiment provides a comprehensive demonstration of
end-to-end machine learning using standard Python libraries.

%%%%%%%%%%%%%%%%%%%%%%%%%%%%%%%%%%%%%%%%%%%%%%%%%%%%%%%%%%%%%%%%%%%%%%
\section{References}
%%%%%%%%%%%%%%%%%%%%%%%%%%%%%%%%%%%%%%%%%%%%%%%%%%%%%%%%%%%%%%%%%%%%%%

\begin{enumerate}

\item F. Pedregosa \textit{et al.},
``Scikit-learn: Machine Learning in Python,''
\textit{Journal of Machine Learning Research},
vol.~12,
pp.~2825--2830,
2011.

\item C. R. Harris \textit{et al.},
``Array Programming with NumPy,''
\textit{Nature},
vol.~585,
pp.~357--362,
2020.

\item W. McKinney,
``Data Structures for Statistical Computing in Python,''
Proceedings of the 9th Python in Science Conference,
pp.~56--61,
2010.

\item P. Virtanen \textit{et al.},
``SciPy 1.0: Fundamental Algorithms for Scientific Computing in Python,''
\textit{Nature Methods},
vol.~17,
pp.~261--272,
2020.

\item J. D. Hunter,
``Matplotlib: A 2D Graphics Environment,''
\textit{Computing in Science \& Engineering},
vol.~9,
no.~3,
pp.~90--95,
2007.

\item Kaggle.
\textit{Iris Species Dataset}.

\url{https://www.kaggle.com/datasets/uciml/iris}

\item Kaggle.
\textit{Pima Indians Diabetes Database}.

\url{https://www.kaggle.com/datasets/uciml/pima-indians-diabetes-database}

\item Kaggle.
\textit{Loan Prediction Problem Dataset}.

\url{https://www.kaggle.com/datasets/altruistdelhite04/loan-prediction-problem-dataset}

\item Kaggle.
\textit{Email Spam Classification Dataset}.

\url{https://www.kaggle.com/datasets/balaka18/email-spam-classification-dataset-csv}

\item Kaggle.
\textit{Handwritten Character Recognition Dataset}.

\url{https://www.kaggle.com/datasets/vaibhao/handwritten-characters}

\end{enumerate}

\end{document}